\section{Discussion}\label{sec:discussion}

\textbf{When does graph learning help?} Our results separate three claims
that are usually conflated. (1) \emph{Graph structure carries signal}: yes.
The TS-KG relations are far more homophilous than the transaction graph
(Table~\ref{tab:homophily}), and causal KG features add
\XgbKgGain{} in F1 to XGBoost
(Table~\ref{tab:main}). The motivation for graph methods on this problem is
therefore real, but it lies in the \emph{entity} structure that a knowledge
graph exposes, not in the transaction-to-transaction graph used by prior GNN
work. (2) \emph{Message passing exploits it}: only partly. On transactions of new
entities, KG-SAGE clearly improves over a graph-free network and over message
passing on the transaction graph, so learning over the KG does capture entity
signal; but it never matches a tree ensemble given the same knowledge, and it
raises many false alarms on recurring licit entities whose neighbourhoods
change over time. High homophily on the training graph is therefore not
sufficient for a GNN advantage under distribution shift, which refines the
conditions conjectured in~\cite{maganti2026}. (3) \emph{Prior reports of
strong performance}: largely a consequence of leaked features, which account for
\XgbLeakDrop{} of the \XgbAllFone{} F1 of XGBoost (\XgbLeakShare{}). For
practitioners, the actionable recommendation is to build the entity knowledge
graph, compute causal features from it, and give them to a well-regularised
tree ensemble with class weights and a validation-tuned threshold; graph
neural networks should be evaluated over several later periods, because a
single validation window can be badly misleading.

\textbf{Label provenance.} Elliptic labels are assigned per controlling
entity of the input addresses~\cite{weber2019,safar2026}. The perfect
homophily of \texttt{CO\_SPEND} is therefore partly a property of the
labelling rule, and the near-perfect homophily of \texttt{CO\_RECEIVE}
suggests that shared output addresses mostly connect transactions of the same
entity (change and deposit addresses). A KG that makes entity structure
explicit thus aligns the model with the structure that defines the labels.
This is also how investigations work in practice, since attributing addresses
to entities is the basis of Bitcoin forensics~\cite{meiklejohn2013}, but it
means that gains measured on Elliptic partly reflect entity re-identification
\emph{within} a time step. Labels produced by a different process, or
entity-disjoint evaluation within a step, may show smaller gains. Because
illicit actors almost never reuse addresses across periods, the
cross-period form of this problem (memorising known entities) is negligible
here, as the unseen-entity view confirms.

\textbf{Threats to validity.} (i) A single dataset, two years of Bitcoin
activity ending in 2017, before SegWit-heavy and Taproot usage. (ii) The
within-step observation assumption: features and edges of a time step are
used to score transactions of that step, i.e.\ scoring happens at the end of
each step; real-time scoring would see fewer neighbours. (iii) The released
Elliptic++ edges cover only transactions inside the benchmark, so address
histories are truncated to the sampled transactions rather than the full
ledger~\cite{safar2026}. (iv) Hyperparameter search was limited by a CPU-only
budget: tree ensembles use standard settings with early stopping, and the
graph models select only the learning rate on the validation window; wider
searches could change the size of the gaps, though not the leakage effect,
which exceeds every other difference we observe. (v) An earlier version of
the \textsc{+Nbr} features encoded transactions without KG neighbours as zero,
which in z-scored space is a valid value; this produced false alarms on
isolated licit transactions and was corrected by encoding them as missing.
All reported numbers use the corrected features.
